\section{Introduction}
\label{sec:introduction}

Expansion planning rests on numbers nobody knows. What fuel will cost a decade
out, how fast demand will grow, and how late projects will actually come online
are all uncertain, and each of them is routinely handled by picking a central
value and running a scenario around it. That practice is defensible only if the
uncertainties are of comparable weight. When one of them dominates, effort spent
hedging the others is effort misallocated, and the planning question changes from
\emph{what might happen} to \emph{which unknown is worth managing}.

Peru is a useful place to ask that question, for three reasons. First, the
generation mix is still built on hydropower and natural gas, with non-conventional
renewables at roughly 5\,\% of generation despite a resource base that modelling
work repeatedly finds economically attractive \citep{fiestas2024,delacruz2024}.
Second, the transmission network is longitudinal: three load zones strung in
series along the coast and the Andes, with the best wind and solar resources
systematically distant from the demand they would serve \citep{valdivia2024}.
Third, the institutional record is mixed in an instructive way. The 1992
restructuring produced genuine competition and improved security of supply
\citep{rudnick2019,perezreyes2009}, and Peru was among the regional pioneers of
technology-specific renewable auctions; but the auction mechanism has since been
suspended, leaving developers without a reliable off-take signal, and new
investment in both generation and transmission slowed once the initial wave of
privatisation ran its course. The country therefore holds, simultaneously, a
strong resource endowment, a grid that was not built for it, and a planning
process whose delivery times are themselves uncertain.

Most of the modelling literature addresses this kind of problem with optimisation.
Generation expansion planning, transmission expansion planning and their joint
formulation are mature fields with a large and still-growing body of work
\citep{gacitua2018,babatunde2020,herding2024,tao2025,valipour2026}. These models
are indispensable for identifying what an efficient expansion path would look
like, but they answer a normative question --- what \emph{should} be built ---
under an assumption of perfect rationality and a single optimising agent. The
system they describe is not the system that exists. In a liberalised market,
generation investment is decentralised and driven by expected profitability,
transmission is planned centrally and built on a different and longer clock, and
the gap between the two is bridged by neither \citep{dyner2001,herrera2019}. It is
precisely in that gap that delays, congestion and stranded capacity appear.

System dynamics was developed for exactly this situation: feedback, accumulation
and delay \citep{forrester1997,sterman2000}. Applied to electricity it has a long
record of reproducing the investment cycles, capacity gaps and price swings that
optimisation abstracts away \citep{ford2001,arango2007,assili2008,zapata2018}, and
recent work continues to use it to evaluate policy in liberalised markets
\citep{luo2024,lin2026,oyarzun2026}. What it is less often asked to do is to
quantify, rather than illustrate, the relative weight of the uncertainties that
drive it. A single simulation --- or a handful of named scenarios --- shows what
the model does; it does not show which input the answer actually depends on.

This paper closes that gap for Peru. We build a system dynamics model of the
Peruvian power system, coupled to a congestion-aware dispatch layer, with an
explicit six-stage regulatory and construction pipeline for both generation and
transmission projects. We then subject it to a designed experiment of 129
simulations: a central case, eight one-at-a-time runs that isolate each driver,
and 120 Latin-hypercube Monte Carlo samples in which four drivers move at once ---
fuel prices, demand growth, and project lead times on the generation and on the
transmission side. Every run covers 2025 to 2050 at a monthly step with an hourly
dispatch inside each month.

The contribution is threefold:

\begin{enumerate}
  \item \textbf{A quantified ranking of the uncertainties, not an illustration of
        them.} We report standardised one-at-a-time sensitivities and rank
        correlations over a Latin-hypercube ensemble, together with a convergence
        check on the reported percentiles. The ranking is unambiguous: demand
        growth dominates every indicator we track, and the other three drivers ---
        including fuel prices, the one most often treated as the central risk for a
        gas-dependent system --- are an order of magnitude weaker.
  \item \textbf{A transmission-explicit representation of a longitudinal system.}
        Expansion is disaggregated into grid built for demand, grid built to
        connect new generation, and grid built to interconnect zones, each with its
        own trigger and its own lead time. This matters in Peru because the three
        categories respond to different drivers and only the third is constrained
        by the series topology.
  \item \textbf{A reproducible artefact.} The model, the input workbook, the
        experimental design, the driver, the analysis code and a commented notebook
        that regenerates every figure are openly available. The validation battery,
        the full parameter set and the extended results are given as supplementary
        material.
\end{enumerate}

The central empirical result is counter-intuitive in a way that bears directly on
policy. Higher demand growth does not raise the long-run average price in this
system; it lowers it, because it pulls forward low-cost renewable capacity that
the central case defers. The least-cost trajectory is consequently also the most
transmission-intensive, requiring about a quarter more network than the central
case. For a system whose grid is already the binding constraint, that is not a
comfortable finding, and it reframes the planning problem: the risk to manage in
Peru is not fuel-price exposure but whether transmission can be delivered fast
enough to serve demand growth that is itself the dominant unknown.

The remainder of the paper is organised as follows. Section~\ref{sec:literature}
reviews the relevant literature on transmission planning under energy transition,
infrastructure delays, renewable integration in emerging economies, and system
dynamics in energy planning, and positions the paper against it.
Section~\ref{sec:peru} describes the Peruvian power system.
Section~\ref{sec:methodology} sets out the model, the experimental design and the
validation protocol. Section~\ref{sec:results} reports the central, best and worst
cases and the uncertainty analysis. Section~\ref{sec:discussion} discusses the
implications and Section~\ref{sec:conclusions} concludes.
